\documentclass[11pt]{article}
\usepackage{cmap}
\usepackage[margin=1.1in]{geometry}
\usepackage{amsmath,amssymb,amsthm}
\usepackage[T1]{fontenc}
\usepackage{lmodern}
\input{glyphtounicode}
\pdfgentounicode=1
\usepackage{xcolor}
\usepackage{hyperref}
\definecolor{linkblue}{rgb}{0.1,0.2,0.55}
\hypersetup{colorlinks=true,linkcolor=linkblue,citecolor=linkblue,urlcolor=linkblue}
\newtheorem{theorem}{Theorem}
\newtheorem{lemma}[theorem]{Lemma}
\newtheorem{corollary}[theorem]{Corollary}
\newtheorem{proposition}[theorem]{Proposition}
\theoremstyle{remark}
\newtheorem{remark}[theorem]{Remark}
\newtheorem{conjecture}[theorem]{Conjecture}

\title{A Wronskian Identity for Antisymmetrized Sine Sums:\\
\large the surface-positivity problem as sine-series ratio
monotonicity, and a counterexample to the naive monotonicity
principle}
\author{Lluis Eriksson}
\date{July 2026}

\begin{document}
\maketitle

\begin{abstract}
For real arrays $(a_m)_{m\ge1}$, $(b_m)_{m\ge1}$ (finitely supported,
or summable with $\sum m(|a_m|+|b_m|) < \infty$) we record the exact
identity
\[
\sum_{1 \le m < n} (a_m b_n - a_n b_m)
\bigl[(m-n)\sin((m+n)t) + (m+n)\sin((n-m)t)\bigr]
\;=\; 2\,\bigl[F_a'(t)\,F_b(t) - F_a(t)\,F_b'(t)\bigr],
\]
where $F_x(t) = \sum_m x_m \sin(mt)$: the antisymmetrized double sum
is a Wronskian of two sine series. The proof is two lines. As a
consequence, the open surface-positivity problem of the companion
$\varphi$-lemma note --- the sign of exactly such a double sum built
from modified-Bessel coefficient arrays --- is \emph{equivalent} to
the strict monotonicity of a ratio of two explicit sine series, and
high-precision numerics indicate that the true statement is global:
the Wronskian is negative on all of $(0,\pi)$, not merely near
$t=\pi$. One might hope this follows from the companion result that
the coefficient ratio $a_m/b_m$ is strictly increasing, in analogy
with the Biernacki--Krzy\.z principle for power series. It does not:
we give a hand-checkable counterexample --- $a = (3,4,3)$,
$b = (3,2,1)$, coefficient ratio $(1,2,3)$ strictly increasing,
$b$ strictly decreasing so that $F_b > 0$ on $(0,\pi)$ by the
classical Fej\'er--Gronwall positivity theorem, yet the Wronskian at
$t = 2\pi/3$ equals $3\sqrt3/2 > 0$ --- so no analogue of
Biernacki--Krzy\.z holds for sine series, and the surface problem
cannot be closed by coefficient-ratio monotonicity alone. We state
the resulting conjecture precisely, record the numerical
phenomenology (near-perfect parity-mirror cancellation of relative
size $e^{-c\beta}$, $c \approx 2.1$, with a working-precision
warning), and identify the structural route we consider most
promising. The identity is machine-checked in Lean~4/Mathlib (finite form,
arbitrary $K$; no \texttt{sorry}, standard axiom oracle), and the
counterexamples are exact algebraic evaluations; nothing beyond the
identity and the counterexamples is claimed.
\end{abstract}

\section{The identity}

\begin{theorem}[Wronskian identity]
\label{thm:wronskian}
Let $(a_m)$, $(b_m)$ be real arrays, finitely supported (or with
$\sum_m m(|a_m| + |b_m|) < \infty$, in which case all series below
converge absolutely and uniformly). Set
$F_a(t) = \sum_{m \ge 1} a_m \sin(mt)$ and likewise $F_b$. Then, for
all $t \in \mathbb{R}$,
\begin{equation}
\label{eq:wronskian}
\sum_{1 \le m < n} (a_m b_n - a_n b_m)
\bigl[(m-n)\sin((m+n)t) + (m+n)\sin((n-m)t)\bigr]
= 2\,\bigl[F_a'(t) F_b(t) - F_a(t) F_b'(t)\bigr].
\end{equation}
\end{theorem}

\begin{proof}
The product-to-sum formulas give, for any $m, n$,
\[
(m-n)\sin((m+n)t) + (m+n)\sin((n-m)t)
= 2\,\bigl[m\cos(mt)\sin(nt) - n\sin(mt)\cos(nt)\bigr] =: 2 G(m,n),
\]
and $G$ is antisymmetric with $G(n,n) = 0$. Since
$(a_m b_n - a_n b_m)$ is antisymmetric as well,
\[
\sum_{m<n} (a_m b_n - a_n b_m)\, 2G(m,n)
= 2 \sum_{m,n} a_m b_n\, G(m,n)
= 2\Bigl[\sum_m a_m m \cos(mt) \sum_n b_n \sin(nt)
- \sum_m a_m \sin(mt) \sum_n b_n n \cos(nt)\Bigr],
\]
which is $2[F_a' F_b - F_a F_b']$.
\end{proof}

\begin{remark}
Equivalently: wherever $F_b \neq 0$, the left side of
\eqref{eq:wronskian} equals $2 F_b^2 \cdot (F_a/F_b)'$. Sign
questions about the antisymmetrized double sum are exactly
monotonicity questions about the ratio of the two sine series ---
with the Wronskian form remaining meaningful (and division-free)
across zeros of $F_b$.
\end{remark}

\section{The surface-positivity problem, reduced}
\label{sec:surface}

We recall the objects of the companion note \cite{PhiLemma} (which
proved the coefficient-ratio monotonicity used below; the wider
programme context is in \cite{MasterMap}). For $\beta > 0$ let
$r_k = I_{k+1}(\beta)/I_1(\beta)$ with $I_\mu$ the modified Bessel
function of the first kind, and for $k \ge 1$
\[
A_k = r_{k-1}^2 \!\!\sum_{j \in \{k-2,\,k\},\, j \ge 0}\!\! (j+1) r_j^2,
\qquad
B_k = k\, r_{k-1}^4,
\qquad
c_{mn} = A_m B_n - A_n B_m .
\]
The arrays satisfy the summability hypothesis of
Theorem~\ref{thm:wronskian} for every $\beta > 0$: by the standard
large-order asymptotics $I_{k}(\beta) \sim (\beta/2)^{k}/k!$
\cite[10.41.1]{DLMF}, the ratios $r_k$ decay super-exponentially in
$k$, hence so do $A_m$ and $B_m$, and
$\sum_m m(A_m + B_m) < \infty$. The open positivity problem is
$W(t) < 0$ for $t$ near $\pi$, where $W$ is the left side of
\eqref{eq:wronskian} with $(a, b) = (A, B)$. By Theorem~\ref{thm:wronskian},
\[
W(t) = 2\,\bigl[F_A'(t) F_B(t) - F_A(t) F_B'(t)\bigr],
\qquad
F_A(t) = \sum_{m\ge1} A_m \sin(mt),\quad
F_B(t) = \sum_{m\ge1} B_m \sin(mt).
\]

Two facts sharpen the picture. First, the companion note proved
$c_{mn} < 0$ for all $m < n$ (the coefficient ratio
$\varphi_m = A_m/B_m$ is strictly increasing) \cite{PhiLemma}.
Second, high-precision numerics (mpmath, $30$--$100$ significant
digits, truncation-stable) indicate that the Wronskian is strictly
negative on \emph{all} of $(0,\pi)$, at every $\beta$ tested
($\beta \in \{0.5, 2, 8, 32\}$, grids of $t/\pi$ from $0.05$ to
$0.99$), and that $F_B > 0$ on $(0,\pi)$. The $\varepsilon$-window
form of the original problem is thus the boundary shadow of a global
statement, which we now record as the target.

\begin{conjecture}[global ratio monotonicity]
\label{conj:main}
For every $\beta > 0$: (i) $F_B > 0$ on $(0,\pi)$, and
(ii) $F_A/F_B$ is strictly decreasing on $(0,\pi)$. Given (i), part
(ii) is equivalent to $W(t) < 0$ on $(0,\pi)$, by
$W = 2F_B^2\,(F_A/F_B)'$; the Wronskian form $W < 0$ is the
division-free statement and is what the numerics test.
\end{conjecture}

\section{The naive principle fails}
\label{sec:counter}

For power series with positive coefficients, the Biernacki--Krzy\.z
principle \cite{BK} states that if $a_m/b_m$ is increasing then
$\sum a_m x^m / \sum b_m x^m$ is increasing on $(0,1)$. Given
$\varphi_m = A_m/B_m$ increasing and $W < 0$ numerically, it is
tempting to conjecture the sine-series analogue: \emph{coefficient
ratio increasing $\Rightarrow$ ratio of sine series decreasing on
$(0,\pi)$} (the sign flip suggested by the boundary behaviour at
$t = \pi$). This is false, even under strong positivity side
conditions.

\begin{proposition}[counterexample, hand-checkable]
\label{prop:counter}
Let $a = (3, 4, 3)$ and $b = (3, 2, 1)$ (all further entries zero).
Then: the coefficient ratio $a_m/b_m = (1, 2, 3)$ is strictly
increasing; $b$ is strictly decreasing, so $F_b(t) = 3\sin t +
2\sin 2t + \sin 3t > 0$ on $(0,\pi)$ by the Fej\'er--Gronwall
positivity theorem for sine polynomials with decreasing coefficients
\cite{AskeySteinig}; and yet at $t = 2\pi/3$,
\[
F_a = -\tfrac{\sqrt3}{2},\quad
F_b = \tfrac{\sqrt3}{2},\quad
F_a' = \tfrac72,\quad
F_b' = -\tfrac12,
\qquad
F_a' F_b - F_a F_b'
= \tfrac{7\sqrt3}{4} - \tfrac{\sqrt3}{4}
= \tfrac{3\sqrt3}{2} \;>\; 0 .
\]
Hence $F_a/F_b$ is strictly \emph{increasing} at $t = 2\pi/3$: no
Biernacki--Krzy\.z analogue holds for sine series.
\end{proposition}

\begin{proof}
At $t = 2\pi/3$: $\sin t = \tfrac{\sqrt3}{2}$,
$\sin 2t = -\tfrac{\sqrt3}{2}$, $\sin 3t = 0$;
$\cos t = -\tfrac12$, $\cos 2t = -\tfrac12$, $\cos 3t = 1$. Then
$F_a = (3 - 4)\tfrac{\sqrt3}{2} = -\tfrac{\sqrt3}{2}$,
$F_b = (3 - 2)\tfrac{\sqrt3}{2} = \tfrac{\sqrt3}{2}$,
$F_a' = 3(-\tfrac12) + 8(-\tfrac12) + 9 = \tfrac72$,
$F_b' = 3(-\tfrac12) + 4(-\tfrac12) + 3 = -\tfrac12$; combine.
\end{proof}

An even smaller example, with the Wronskian form: $a = (1,2,3,4)$,
$b = (1,1,1,1)$ gives $F_a'F_b - F_aF_b' = 4$ at $t = \pi/2$
(there $F_b$ vanishes, so this violates $W \le 0$ though not
ratio-monotonicity on a domain of positivity --- which is why
Proposition~\ref{prop:counter} carries the decreasing-$b$
condition).

\begin{remark}[the moral, in the spirit of the parity barrier]
\label{rem:moral}
As with the parity barrier for decoupling inequalities
\cite{ParityBarriers}, a bounded amount of structural data ---
here, the monotonicity of the coefficient ratio, even together with
positivity of $F_b$ --- cannot certify the global sign. Whatever
proves Conjecture~\ref{conj:main} must consume more of the Bessel
structure of $(A_m, B_m)$ than the order of $\varphi_m$. This is
consistent with the quantitative phenomenology: the even- and
odd-parity halves of $W(\pi - \varepsilon)$ cancel to relative
accuracy $\approx e^{-c\beta}$ with $c \approx 2.03$--$2.17$ (true
totals e.g.\ $-2.9\cdot10^{-8}$ at $\beta=8$;
$-5.0\cdot10^{-29}$ at $\beta=32$; $-5.1\cdot10^{-95}$ at
$\beta=100$), so any estimate that discards the exact structure
loses the sign entirely; numerical exploration below
$\approx 2.2\beta + 20$ working digits sees pure rounding noise.
The route we consider most promising: $r_{k}^2$ and $r_k^4$ are
ratios of Bessel products with integral (Schl\"afli/heat-kernel)
representations, so $F_A$ and $F_B$ admit kernel representations on
the circle; a positive integral representation of
$-(F_A/F_B)'$ --- or a total-positivity property of the array
$(A; B)$ against the sine kernel --- would close
Conjecture~\ref{conj:main} where coefficient-order arguments
provably cannot.
\end{remark}

\section{Verification}

Theorem~\ref{thm:wronskian} was verified symbolically (sympy) for
generic coefficients at $K \le 5$ (residual $0$ after trigonometric
normalisation) and at $K = 7$ by exact rational substitution
(residual $< 10^{-35}$ at $40$ digits, four random points), and
numerically against the direct double sum for the Bessel arrays at
$K = 90$, $\beta = 8$, at $t \in \{0.5, 1.5, 2.6, \pi - 0.1\}$
(ratio $1.0$ to all displayed digits). The counterexample values of
Proposition~\ref{prop:counter} are exact algebraic evaluations. The
global sign scan of $W$ and the parity-mirror measurements are as in
the companion note. Scripts and sources are archived at
\begin{center}\small
\url{https://github.com/lluiseriksson/THE-ERIKSSON-PROGRAMME}
(\texttt{docs/BF2-ATTACK-NOTES.md}, \texttt{papers/wronskian-reduction}).
\end{center}
Theorem~\ref{thm:wronskian} (finite form, arbitrary $K$) is
moreover machine-checked in Lean~4/Mathlib (file
\texttt{WronskianIdentity.lean} accompanying this note; checked with
\texttt{lake env lean} against Lean \texttt{4.30.0-rc2}, Mathlib
commit \texttt{cd3b69b}): \texttt{bracket\_eq\_two\_G} (the
trigonometric collapse), \texttt{double\_sum\_eq\_wronskian} (the
product structure), \texttt{pair\_sum\_eq\_double\_sum} (the
antisymmetrization) and \texttt{wronskian\_identity} (the theorem).
The axiom oracle returns Lean's three standard axioms for all four;
no \texttt{sorry}:
\begin{verbatim}
'wronskian_identity' depends on axioms:
    [propext, Classical.choice, Quot.sound]
\end{verbatim}
(exit code 0; likewise for the three lemmas).

\subsection*{Acknowledgements}
This note is part of an audit-first research programme whose records,
Lean sources and numerical suites are public at
\url{https://github.com/lluiseriksson/THE-ERIKSSON-PROGRAMME}.
Nothing is claimed about Conjecture~\ref{conj:main} beyond its
equivalence to the original surface-positivity problem and the
failure of the naive route; nothing is claimed about
four-dimensional gauge theory.

\begin{thebibliography}{9}\small

\bibitem{BK}
M.~Biernacki and J.~Krzy\.z,
\emph{On the monotonity of certain functionals in the theory of
analytic functions},
Ann. Univ. Mariae Curie-Sk{\l}odowska Sect. A \textbf{9} (1955),
135--147.

\bibitem{AskeySteinig}
R.~Askey and J.~Steinig,
\emph{Some positive trigonometric sums},
Trans. Amer. Math. Soc. \textbf{187} (1974), 295--307.

\bibitem{DLMF}
F.~W.~J. Olver et al.,
\emph{NIST Digital Library of Mathematical Functions},
\S10.29, \url{https://dlmf.nist.gov/10.29}.

\bibitem{PhiLemma}
L.~Eriksson,
\emph{A weighted Tur\'an-type monotonicity lemma for modified Bessel
functions via the calibrated Amos bound, with an application to the
ordering of a surface expansion in two-dimensional lattice gauge
theory},
July 2026;
\url{https://github.com/lluiseriksson/THE-ERIKSSON-PROGRAMME/tree/main/papers/phi-lemma}.

\bibitem{ParityBarriers}
L.~Eriksson,
\emph{Parity barriers for decoupling inequalities},
July 2026;
\url{https://github.com/lluiseriksson/THE-ERIKSSON-PROGRAMME/tree/main/papers/parity-barriers}.

\bibitem{MasterMap}
L.~Eriksson,
\emph{THE MASTER MAP: an audit-first navigation guide to the
conditional construction of 4D $SU(N)$ Yang--Mills with mass gap},
ai.viXra:2602.0096, and the programme repository
\url{https://github.com/lluiseriksson/THE-ERIKSSON-PROGRAMME}.

\end{thebibliography}

\end{document}
